\documentclass[11pt, a4paper]{article}
\usepackage{fontspec}
\usepackage{xeCJK}
\setCJKmainfont{Microsoft JhengHei}
\usepackage{geometry}
\geometry{left=1.8cm, right=1.8cm, top=2cm, bottom=2cm}
\usepackage{xcolor}
\usepackage{colortbl}
\usepackage{hyperref}
\usepackage{tcolorbox}
\usepackage{booktabs}
\usepackage{array}

\definecolor{traBlue}{RGB}{24, 76, 120}
\definecolor{traDarkRed}{RGB}{168, 50, 50}
\definecolor{traGreen}{RGB}{34, 139, 34}
\definecolor{lightGray}{RGB}{245, 247, 250}

\hypersetup{
    colorlinks=true,
    linkcolor=traBlue,
    urlcolor=traBlue
}

\title{\textbf{\Huge TR-PASS 學生版 5日全台鐵道遊計劃手冊}\\[0.3em]
\Large -- 板橋起訖・嚴格合規車次・轉乘矩陣與備案指南 --}
\author{\textbf{台鐵 2026/07/01 最新時刻表大數據規劃}}
\date{\today}

\begin{document}

\maketitle

\begin{tcolorbox}[colback=traBlue!10!white,colframe=traBlue,title=\textbf{TR-PASS 學生版乘車金科玉律 (法規核驗 100\% 合格)}]
\begin{enumerate}
    \item \textbf{絕對禁搭 (違者補票並加收 50\%)}：普悠瑪號、太魯閣號、\textbf{EMU3000 型新自強號}、觀光/團體專列。
    \item \textbf{完全適用 (全免劃位，刷 Pass 進站)}：區間車、\textbf{區間快車 (EMU900/EMU800)}、莒光號、\textbf{一般 PP 推挽式自強號 (自由席無座)}。
    \item \textbf{每日核心守則}：每天皆為 \textbf{「板橋站早出 $\to$ 當晚返抵板橋站」}。
\end{enumerate}
\end{tcolorbox}

\section{5日行程總覽與時序對照}

\begin{table}[h!]
\centering
\rowcolors{2}{lightGray}{white}
\begin{tabular}{p{1.2cm}p{1.8cm}p{3.8cm}p{3.8cm}p{1.5cm}p{3.8cm}}
\toprule
\textbf{天數} & \textbf{目的地} & \textbf{去程班次 (板橋出發)} & \textbf{回程班次 (返抵板橋)} & \textbf{停留時間} & \textbf{主要景點與特色} \\
\midrule
\textbf{Day 1} & 台中/彰化 & 區快 2011 (07:40 $\to$ 09:50) & 區快 2036 (17:40 $\to$ 19:55) & 7.5 小時 & 台中舊站、第二市場、彰化扇形車庫 \\
\textbf{Day 2} & 台東/花蓮 & 區快 2005+莒光 727 (06:15 $\to$ 12:40) & 區快 4030 (18:50 $\to$ 21:55) & 2.75 小時 & 南迴太平洋海景、傳奇莒光 554 長征 \\
\textbf{Day 3} & 平溪/內灣 & 區間 4148 (08:10 $\to$ 09:05) & 區快 2038 (18:30 $\to$ 19:40) & 7.6 小時 & 八斗子無敵海景、十分天燈、內灣老街 \\
\textbf{Day 4} & 集集/成追 & PP自強 109 (07:46 $\to$ 10:10) & 區快 2036 (17:40 $\to$ 19:55) & 7.0 小時 & 集集綠色隧道、車埕便當、追分成功車票 \\
\textbf{Day 5} & 花東縱谷 & 區快 4003+4022 (06:50 $\to$ 12:25) & 自強 177 (20:00 $\to$ 23:05) & 5.5 小時 & 池上單車伯朗大道、花蓮東大門夜市 \\
\bottomrule
\end{tabular}
\end{table}

\section{每日詳細車次、轉乘與 Plan B 備案}

\subsection{Day 1：西部幹線區間快極速遊（板橋 $\leftrightarrow$ 台中 / 彰化）}
\begin{itemize}
    \item \textbf{07:40 - 09:50} ｜ \textbf{【去程】} 搭乘 \textbf{區間快 2011次 (EMU900)}：板橋 07:40 $\to$ 台中 09:50。（車上享用早餐）
    \item \textbf{10:00 - 13:00} ｜ \textbf{【台中】} 遊覽台中車站舊站區、綠川水岸，於 \textbf{第二市場享用午餐}（王家菜頭粿、山河爌肉飯）。
    \item \textbf{13:08 - 13:20} ｜ \textbf{【轉乘】} 搭乘 \textbf{區間車 2173次}：台中 13:08 $\to$ 彰化 13:20。
    \item \textbf{13:30 - 17:00} ｜ \textbf{【彰化】} 參觀國定古蹟 \textbf{彰化扇形車庫}，品嚐 \textbf{阿三肉圓}。
    \item \textbf{17:40 - 19:55} ｜ \textbf{【回程】} 搭乘 \textbf{區間快 2036次 (EMU900)}：彰化 17:40 $\to$ 板橋 19:55 返抵。
\end{itemize}
\begin{tcolorbox}[colback=yellow!10!white,colframe=orange,title=\textbf{Day 1 應變備案 (Plan B)}]
若早上錯過 07:40 區間快，可改搭 \textbf{PP自強號 109次}（板橋 07:46 發車 $\to$ 台中 09:55 抵達），時間完全無縫接軌！
\end{tcolorbox}

\subsection{Day 2：南迴海景與莒光號長征（板橋 $\leftrightarrow$ 枋寮 $\leftrightarrow$ 花蓮 $\leftrightarrow$ 板橋）}
\begin{itemize}
    \item \textbf{06:15 - 10:28} ｜ 板橋 06:15 $\to$ 新左營 09:40 (區間快2005) $\to$ 枋寮 10:28 (區間快3063)。
    \item \textbf{10:45 - 12:40} ｜ 枋寮 10:45 $\to$ 台東 12:40 (搭乘 \textbf{莒光號 727次}，遠眺多良太平洋海景)。
    \item \textbf{12:57 - 17:30} ｜ 體驗全台對號最慢車 \textbf{莒光號 554次}：台東 12:57 $\to$ 花蓮 17:30（車上享用池上便當）。
    \item \textbf{17:30 - 18:45} ｜ \textbf{【花蓮站前停留 75 分鐘】} 步行品嚐 \textbf{液香扁食 / 炸蛋蔥油餅}。
    \item \textbf{18:50 - 21:55} ｜ 搭乘 \textbf{區間快 4030次 (EMU900)}：花蓮 18:50 $\to$ 板橋 21:55 返抵。
\end{itemize}

\subsection{Day 3：北部雙支線慢活（平溪/深澳線 + 內灣線）}
\begin{itemize}
    \item \textbf{08:10 - 09:05} ｜ 板橋 08:10 $\to$ 瑞芳 09:05 (區間 4148次)。
    \item \textbf{09:25 - 12:30} ｜ 平溪深澳線：遊覽八斗子海景月台與十分老街（享用老街小吃）。
    \item \textbf{13:00 - 14:35} ｜ 瑞芳 13:00 $\to$ 新竹 14:35 (區間快 2017次)。
    \item \textbf{14:45 - 18:00} ｜ 內灣線：漫步內灣老街（品嚐野薑花粽、客家擂茶）。
    \item \textbf{18:30 - 19:40} ｜ 搭乘 \textbf{區間快 2038次 (EMU900)}：新竹 18:30 $\to$ 板橋 19:40 返抵。
\end{itemize}

\subsection{Day 4：中部懷舊與成追線（集集線 + 追分成功）}
\begin{itemize}
    \item \textbf{07:46 - 10:10} ｜ 搭乘 \textbf{PP自強號 109次}：板橋 07:46 $\to$ 二水 10:10 抵達。
    \item \textbf{10:30 - 14:00} ｜ 集集線：二水 10:30 $\to$ 車埕 11:20（車埕木業園區享用 \textbf{車埕木桶便當}）。
    \item \textbf{14:30 - 16:30} ｜ 前往成功/追分站，購買「追分成功」紀念車票。
    \item \textbf{17:40 - 19:55} ｜ 搭乘 \textbf{區間快 2036次}：彰化 17:40 $\to$ 板橋 19:55 返抵。
\end{itemize}

\subsection{Day 5：東幹線區間快大縱走（板橋 $\leftrightarrow$ 池上單車 $\leftrightarrow$ 花蓮夜市）}
\begin{itemize}
    \item \textbf{06:50 - 09:40} ｜ 東線區間快 4003次：板橋 06:50 $\to$ 花蓮 09:40。
    \item \textbf{10:00 - 12:25} ｜ 縱谷區間快 4022次：花蓮 10:00 $\to$ 池上 12:25。
    \item \textbf{12:30 - 15:50} ｜ \textbf{【池上】} 吃 \textbf{池上悟饕木盒便當}，租單車暢遊伯朗大道、天堂路。
    \item \textbf{16:00 - 18:25} ｜ 區間車 4537次：池上 16:00 $\to$ 花蓮 18:25。
    \item \textbf{18:30 - 20:00} ｜ \textbf{【花蓮停留 90 分鐘】} 逛 \textbf{東大門夜市} 或品嚐公正包子。
    \item \textbf{20:00 - 23:05} ｜ 搭乘 \textbf{PP自強號 177次}：花蓮 20:00 $\to$ 板橋 23:05 返抵。
\end{itemize}

\end{document}
